% 01_overview.tex - Project Overview
% AlohaMini Codebase Documentation

\section{Project Overview}

% ============================================
% SLIDE: Introduction
% ============================================
\begin{frame}{AlohaMini Robot}
\begin{columns}
\begin{column}{0.5\textwidth}
    \textbf{Dual-Arm Mobile Manipulation Robot}
    \vspace{0.5cm}

    \begin{itemize}
        \item 16 DOF total configuration
        \item Omnidirectional mobile base (3 wheels)
        \item Vertical lift mechanism
        \item Dual 6-DOF arms
        \item 5-camera perception system
    \end{itemize}
\end{column}
\begin{column}{0.5\textwidth}
    \centering
    \begin{tikzpicture}[scale=0.8, transform shape]
        % Base
        \draw[fill=gray!30, rounded corners=5pt] (-2,-0.5) rectangle (2,0.5);
        \node at (0,0) {\small Mobile Base};

        % Wheels
        \draw[fill=black!60] (-1.5,-0.7) circle (0.2);
        \draw[fill=black!60] (1.5,-0.7) circle (0.2);
        \draw[fill=black!60] (0,-0.7) circle (0.2);

        % Lift
        \draw[fill=blue!20] (-0.3,0.5) rectangle (0.3,2);
        \node[rotate=90] at (0,1.25) {\tiny Lift};

        % Arm platform
        \draw[fill=gray!20] (-1.5,2) rectangle (1.5,2.5);

        % Left arm
        \draw[line width=2pt, gray!60] (-1,2.5) -- (-1.5,3.5) -- (-2,4);
        \node[joint] at (-1,2.5) {};
        \node[joint] at (-1.5,3.5) {};
        \node[joint] at (-2,4) {};
        \node[above] at (-2,4) {\tiny L};

        % Right arm
        \draw[line width=2pt, gray!60] (1,2.5) -- (1.5,3.5) -- (2,4);
        \node[joint] at (1,2.5) {};
        \node[joint] at (1.5,3.5) {};
        \node[joint] at (2,4) {};
        \node[above] at (2,4) {\tiny R};
    \end{tikzpicture}
\end{column}
\end{columns}
\end{frame}

% ============================================
% SLIDE: Robot Specifications
% ============================================
\begin{frame}{Robot Specifications}
\begin{table}
\centering
\begin{tabular}{lcc}
\toprule
\textbf{Component} & \textbf{DOF} & \textbf{Type} \\
\midrule
Mobile Base (X, Y, $\theta$) & 3 & Prismatic/Continuous \\
Vertical Lift & 1 & Prismatic \\
Left Arm & 6 & Revolute \\
Right Arm & 6 & Revolute \\
\midrule
\textbf{Total} & \textbf{16} & \\
\bottomrule
\end{tabular}
\end{table}

\vspace{0.5cm}

\begin{block}{Action Space (ManiSkill3)}
\texttt{[base\_x, base\_y, base\_rot, lift, left\_arm(6), right\_arm(6)]}
\end{block}
\end{frame}

% ============================================
% SLIDE: Software Stack
% ============================================
\begin{frame}{Software Stack}
\centering
\begin{tikzpicture}[node distance=0.8cm, scale=0.9, transform shape]
    % Stack layers
    \node[module, text width=8cm, minimum height=1cm] (app) at (0,0)
        {Application Layer\\(VR Teleop, Keyboard Control, Voice Commands)};

    \node[module, text width=8cm, minimum height=1cm, below=of app] (ctrl)
        {Control Layer\\(TeleopController, Kinematics, Input Providers)};

    \node[module, text width=8cm, minimum height=1cm, below=of ctrl] (sim)
        {Simulation Layer\\(ManiSkill3, SAPIEN, PyTorch)};

    \node[module, text width=8cm, minimum height=1cm, below=of sim] (hw)
        {Hardware Layer\\(LeRobot, Feetech Motors, Cameras)};

    % Arrows
    \draw[flow] (app) -- (ctrl);
    \draw[flow] (ctrl) -- (sim);
    \draw[flow] (ctrl) -- (hw);

    % Side labels
    \node[left=0.5cm of app, font=\footnotesize, text=gray] {Python/WebXR};
    \node[left=0.5cm of ctrl, font=\footnotesize, text=gray] {teleop/};
    \node[left=0.5cm of sim, font=\footnotesize, text=gray] {agents/};
    \node[left=0.5cm of hw, font=\footnotesize, text=gray] {software/};
\end{tikzpicture}
\end{frame}

% ============================================
% SLIDE: Key Features
% ============================================
\begin{frame}{Key Features}
\begin{columns}[t]
\begin{column}{0.5\textwidth}
    \textbf{Simulation (ManiSkill3)}
    \begin{itemize}
        \item GPU-accelerated physics
        \item Virtual mobile base
        \item Grasping detection
        \item Camera rendering
        \item ReplicaCAD scenes
    \end{itemize}

    \vspace{0.3cm}

    \textbf{Teleoperation}
    \begin{itemize}
        \item Keyboard control (XLeRobot-style)
        \item VR control (WebXR)
        \item 50 Hz control loop
        \item IK-based arm control
    \end{itemize}
\end{column}
\begin{column}{0.5\textwidth}
    \textbf{Physical Robot}
    \begin{itemize}
        \item LeRobot integration
        \item 5-camera perception
        \item Trajectory recording
        \item Voice commands
    \end{itemize}

    \vspace{0.3cm}

    \textbf{ROS Simulation}
    \begin{itemize}
        \item RViz visualization
        \item Gazebo physics
        \item URDF models
        \item ROS1/ROS2 support
    \end{itemize}
\end{column}
\end{columns}
\end{frame}

% ============================================
% SLIDE: Code Metrics
% ============================================
\begin{frame}{Code Metrics}
\begin{columns}
\begin{column}{0.5\textwidth}
    \begin{table}
    \centering
    \begin{tabular}{lr}
    \toprule
    \textbf{Module} & \textbf{Lines} \\
    \midrule
    agents/aloha\_mini & $\sim$850 \\
    teleop/ & $\sim$2,000 \\
    kinematics/ & $\sim$800 \\
    software/ (LeRobot) & $\sim$1,500 \\
    \midrule
    \textbf{Total} & $\sim$5,150 \\
    \bottomrule
    \end{tabular}
    \end{table}
\end{column}
\begin{column}{0.5\textwidth}
    \textbf{Key Classes}
    \begin{itemize}
        \item \code{AlohaMiniVirtual}
        \item \code{TeleopController}
        \item \code{AlohaMiniKinematics}
        \item \code{ControlGoal}
        \item \code{KeyboardController}
        \item \code{VRControllerServer}
    \end{itemize}
\end{column}
\end{columns}
\end{frame}
